\chapter[Hierarchical VLA and Embodied Reasoning]{Hierarchical VLA and\\Embodied Reasoning}

\section*{학습 목표}

이 장은 Real VLA를 단일 거대 policy가 아니라 여러 시간 척도와 책임 경계를 가진 계층형 system으로 본다. 상위 layer는 장면과 명령을 해석하고, task를 subgoal로 나누고, 진행 상황과 실패를 판단한다. 하위 VLA policy는 주어진 subgoal 또는 instruction에 대해 action chunk, waypoint, trajectory를 생성한다. controller와 safety monitor는 물리적 실행을 안정화하고 제한한다.

\begin{enumerate}
  \item embodied reasoning layer의 역할과 출력 형식을 정의한다.
  \item hierarchical policy를 options와 belief-state 관점에서 정식화한다.
  \item reasoning layer, VLA policy, controller, safety monitor 사이의 interface contract를 설계한다.
  \item progress detection, success estimation, replanning trigger를 평가 가능한 형태로 만든다.
\end{enumerate}

\section{Why hierarchy returns}

초기 VLA 논의는 end-to-end policy의 가능성을 강조했다. 관측과 언어를 넣으면 action이 직접 나오는 구조는 data가 충분하고 task horizon이 짧을 때 강력하다 \cite{rt12022,openvla2024,pi02024}.
\begin{equation}
  \act_t
  \sim
  \pi_{\theta}
  (\cdot \mid \obs_{\leq t}, \ell).
  \label{eq:ch17-monolithic-policy}
\end{equation}
그러나 실제 배치에 가까운 task에서는 이 단순한 interface만으로 부족한 경우가 많다. 이유는 네 가지다.
\begin{enumerate}
  \item task horizon이 길어질수록 action error보다 subgoal error가 더 중요해진다.
  \item safety constraint는 policy 내부의 implicit preference가 아니라 외부에서 검증 가능한 조건이어야 한다.
  \item reasoning과 action generation은 요구하는 시간 척도와 compute budget이 다르다.
  \item 실패 후 복구는 단순 next action prediction이 아니라 상태 재해석과 계획 수정이 필요하다.
\end{enumerate}

따라서 real deployment에서는 다음 구조가 자연스럽다.
\begin{equation}
  z_t
  =
  \rho_{\psi}
  (\ell, \obs_{\leq t}, \bel_t, m_t),
  \qquad
  \act_{t:t+H-1}
  \sim
  \pi_{\theta}
  (\cdot \mid \obs_{\leq t}, z_t, \conf_t).
  \label{eq:ch17-reasoning-vla}
\end{equation}
여기서 $\rho_{\psi}$는 embodied reasoning layer, $z_t$는 subgoal, constraint, success criterion, tool call, recovery instruction을 포함하는 중간 표현이다. 하위 VLA는 $z_t$를 실행 가능한 action sequence로 변환한다.

\begin{definitionbox}{Hierarchical VLA}
Hierarchical VLA는 language-conditioned reasoning layer와 action-generating VLA policy를 분리하되, 두 layer가 동일한 관측, task context, progress state를 공유하도록 설계한 robot policy stack이다. 핵심은 모듈을 많이 나누는 것이 아니라 각 layer의 책임과 출력 검증 조건을 명확히 하는 것이다.
\end{definitionbox}

\section{Embodied reasoning is not text reasoning}

LLM의 text reasoning은 문장 사이의 논리 관계를 다룬다. embodied reasoning은 물리 세계 안에서 실행될 action의 조건을 다룬다. 같은 명령이라도 robot embodiment, object pose, gripper state, scene constraint에 따라 의미가 달라진다.

명령 $\ell$이 주어졌을 때 embodied reasoning layer가 추정해야 할 변수는 다음과 같이 쓸 수 있다.
\begin{equation}
  r_t
  =
  (
  g_t,
  c_t,
  \sigma_t,
  q_t,
  \eta_t
  ).
  \label{eq:ch17-reasoning-state}
\end{equation}
여기서 $g_t$는 현재 subgoal, $c_t$는 constraint set, $\sigma_t$는 success criterion, $q_t$는 필요한 clarification 또는 perception query, $\eta_t$는 recovery mode이다.

\begin{table}[h]
  \centering
  \caption{Embodied reasoning의 주요 출력}
  \begin{tabularx}{\linewidth}{p{0.22\linewidth}X X}
    \toprule
    Output & 의미 & 예시 \\
    \midrule
    Subgoal $g_t$ & 현재 실행해야 할 물리 목표 & 컵을 오른쪽 gripper workspace 안으로 이동 \\
    Constraint $c_t$ & 실행 중 지켜야 할 조건 & 칼날 방향 회피, 사람 근처 속도 제한 \\
    Success criterion $\sigma_t$ & subgoal 완료 판정 조건 & object centroid가 bin 내부이고 gripper가 열림 \\
    Query $q_t$ & 추가 관측 또는 사람 질문 & 가려진 물체 확인, ambiguous instruction 확인 \\
    Recovery mode $\eta_t$ & 실패 후 전환할 행동 범주 & regrasp, re-detect, ask, abort \\
    \bottomrule
  \end{tabularx}
  \label{tab:ch17-er-output}
\end{table}

중요한 점은 reasoning output이 자연어 설명으로 끝나면 안 된다는 것이다. 하위 VLA와 controller가 사용할 수 있는 contract로 내려와야 한다. 예를 들어 ``파란 컵을 집어라''는 문장은 robot stack 내부에서 object id, target pose distribution, allowed approach direction, gripper command convention, success predicate로 변환되어야 한다.

\section{Options view of hierarchical VLA}

계층형 VLA는 options framework로 정식화할 수 있다. 상위 policy는 option $\omega_t$를 선택하고, 하위 policy는 선택된 option 아래에서 action을 생성한다.
\begin{equation}
  \omega_t
  \sim
  \mu_{\psi}
  (\cdot \mid \bel_t, \ell),
  \qquad
  \act_t
  \sim
  \pi_{\theta}
  (\cdot \mid \obs_t, \omega_t).
  \label{eq:ch17-option-policy}
\end{equation}
option은 단순 skill 이름일 수도 있고, continuous subgoal, constraint set, language instruction을 함께 가진 구조화된 object일 수도 있다.

각 option은 termination model을 갖는다.
\begin{equation}
  \beta_{\omega}(\bel_t)
  =
  P(\mathrm{terminate}=1 \mid \bel_t,\omega).
  \label{eq:ch17-option-termination}
\end{equation}
termination은 성공뿐 아니라 실패, unsafe condition, ambiguity, timeout도 포함한다. 따라서 hierarchical VLA의 상위 loop는 다음과 같이 쓸 수 있다.
\begin{equation}
  \omega_{t+1}
  =
  \begin{cases}
  \mu_{\psi}(\bel_t,\ell) & \text{if } \beta_{\omega_t}(\bel_t)=1, \\
  \omega_t & \text{otherwise.}
  \end{cases}
  \label{eq:ch17-option-loop}
\end{equation}

이 관점은 VLA를 classical skill library로 되돌리는 것이 아니다. option 내부의 하위 policy가 여전히 visuomotor foundation model일 수 있고, option 자체도 language-conditioned reasoning model이 생성할 수 있다. 차이는 action prediction을 무한히 길게 늘리는 대신, 긴 horizon을 검증 가능한 subproblem으로 나누는 데 있다.

\section{Interface contract between reasoning and VLA}

계층 구조가 잘 작동하려면 layer 사이의 interface가 느슨한 자연어 prompt가 아니라 typed contract에 가까워야 한다. 한 time step에서 reasoning layer가 VLA에 전달하는 command를 다음 tuple로 정의하자.
\begin{equation}
  z_t
  =
  (
  \ell_t^{sub},
  \mathcal{O}_t,
  \mathcal{G}_t,
  \mathcal{C}_t,
  \Sigma_t,
  H_t
  ).
  \label{eq:ch17-command-contract}
\end{equation}
여기서 $\ell_t^{sub}$는 sub-instruction, $\mathcal{O}_t$는 relevant object set, $\mathcal{G}_t$는 target geometry 또는 pose distribution, $\mathcal{C}_t$는 constraint, $\Sigma_t$는 success predicate, $H_t$는 action horizon이다.

\begin{table}[h]
  \centering
  \caption{Reasoning--VLA interface contract}
  \begin{tabularx}{\linewidth}{p{0.22\linewidth}X X}
    \toprule
    Field & 필요 이유 & 검증 방법 \\
    \midrule
    Sub-instruction & 하위 policy의 local objective 지정 & instruction grounding check \\
    Object set & attention과 affordance 범위 제한 & detector id, mask, pose confidence \\
    Geometry & motion target 또는 relation 지정 & frame convention, pose covariance \\
    Constraint & collision, force, human proximity 제한 & safety monitor predicate \\
    Success predicate & option termination 조건 제공 & state estimator 또는 verifier \\
    Horizon & action chunk 길이와 interruptibility 결정 & latency budget, override test \\
    \bottomrule
  \end{tabularx}
  \label{tab:ch17-contract}
\end{table}

contract가 없으면 상위 model은 그럴듯한 subgoal을 말하고, 하위 VLA는 그 subgoal을 다르게 해석하며, controller는 불가능한 target을 받는다. 이 mismatch는 성능 문제가 아니라 system integration 문제다.

\section{Progress detection and success estimation}

long-horizon task에서 핵심은 다음 action보다 progress state이다. progress를 scalar로만 두면 불충분하다. subgoal별 completion, constraint violation, uncertainty, recovery need를 함께 추정해야 한다.
\begin{equation}
  p_t
  =
  \chi_{\xi}
  (\obs_{\leq t}, z_t, \tau_{0:t})
  =
  (
  P_{done},
  P_{fail},
  P_{unsafe},
  U_{state},
  U_{intent}
  ).
  \label{eq:ch17-progress-state}
\end{equation}
여기서 $\chi_{\xi}$는 progress detector 또는 verifier이다. $U_{state}$는 state uncertainty, $U_{intent}$는 instruction ambiguity를 나타낸다.

replanning trigger는 다음 조건으로 정의할 수 있다.
\begin{equation}
  \mathrm{replan}_t
  =
  \mathbb{I}
  [
  P_{done} > \alpha
  \vee
  P_{fail} > \gamma
  \vee
  P_{unsafe} > \delta
  \vee
  U_{state} > \epsilon
  \vee
  t-t_0 > T_{\max}
  ].
  \label{eq:ch17-replan-trigger}
\end{equation}
이 식은 중요한 engineering principle을 말한다. reasoning은 매 step마다 반드시 호출될 필요가 없다. 그러나 완료, 실패, 위험, 불확실성, timeout이 감지되면 상위 layer가 다시 개입해야 한다.

\section{VLA plus TAMP}

Task and Motion Planning은 symbolic precondition과 geometric feasibility를 함께 다룬다. VLA와 TAMP를 결합할 때 TAMP는 반드시 모든 action을 결정하는 central planner일 필요가 없다. 세 가지 역할이 가능하다.

\begin{enumerate}
  \item 상위 reasoning이 symbolic subgoal sequence를 만들고, TAMP가 feasibility를 검증한다.
  \item TAMP가 candidate plan skeleton을 만들고, VLA가 각 step의 visual execution을 담당한다.
  \item VLA가 proposed action을 만들고, TAMP 또는 geometric checker가 constraint violation을 걸러낸다.
\end{enumerate}

이를 optimization 형태로 쓰면 다음과 같다.
\begin{equation}
  \max_{\omega_{0:K}, \act_{0:T}}
  \mathbb{E}
  \left[
  \sum_{t=0}^{T} r(\state_t,\act_t,\ell)
  -
  \lambda
  \sum_{t=0}^{T}
  \mathbb{I}[\state_t \notin \mathcal{S}_{safe}]
  \right]
  \quad
  \mathrm{s.t.}\quad
  \act_t \sim \pi_{\theta}(\cdot \mid \obs_t,\omega_k).
  \label{eq:ch17-hier-objective}
\end{equation}
여기서 $\omega_{0:K}$는 subgoal 또는 option sequence이고, $\mathcal{S}_{safe}$는 safety set이다. TAMP의 가치는 reward를 크게 만드는 것보다 infeasible branch를 조기에 제거하는 데서 더 크게 나타날 수 있다.

\section{Tool and function calling}

embodied reasoning layer는 종종 외부 tool을 호출해야 한다. 예를 들어 object detector, pose estimator, mapping module, grasp sampler, collision checker, memory retrieval, human query interface가 tool이 될 수 있다.
\begin{equation}
  y_t^{tool}
  =
  \mathcal{T}_{k_t}
  (x_t),
  \qquad
  k_t
  \sim
  \rho_{\psi}
  (\cdot \mid \ell,\obs_{\leq t},\bel_t).
  \label{eq:ch17-tool-call}
\end{equation}
tool calling은 LLM application에서 흔하지만, robot에서는 tool output이 physical action의 전제 조건이 된다. 따라서 tool result에는 confidence, frame convention, timestamp, validity region이 포함되어야 한다.

\begin{warningbox}{Tool call의 실패 모드}
tool call이 성공했다는 software-level return code와 robot task가 성공할 조건은 다르다. pose estimator가 값을 반환했더라도 frame이 틀리거나 오래된 observation이면 하위 VLA는 잘못된 target을 실행한다. 따라서 tool result는 controller가 쓰기 전에 freshness와 consistency를 검사해야 한다.
\end{warningbox}

\section{World models and memory}

hierarchical VLA에서 world model은 두 수준에서 쓰인다. 첫째, short-horizon action consequence를 예측하여 unsafe 또는 infeasible action을 걸러낸다. 둘째, long-horizon task state를 유지하여 이미 완료된 subgoal과 남은 subgoal을 구분한다.
\begin{equation}
  \hat{\state}_{t+1}
  =
  f_{\varphi}
  (\bel_t,\act_t),
  \qquad
  m_{t+1}
  =
  \mathcal{M}(m_t,\obs_t,\act_t,\hat{\state}_{t+1}).
  \label{eq:ch17-world-memory}
\end{equation}
여기서 $m_t$는 scene memory 또는 episodic memory이다. memory가 없으면 robot은 같은 drawer를 반복해서 열거나 이미 옮긴 object를 다시 찾을 수 있다.

world model은 완벽할 필요가 없다. 실용적으로는 다음 세 질문에 답하면 충분한 경우가 많다.
\begin{enumerate}
  \item 이 action이 가까운 미래에 collision 또는 joint limit violation을 만들 가능성이 큰가?
  \item 현재 subgoal이 이미 완료되었는가?
  \item 실패했을 때 이전에 성공한 recovery pattern이 있는가?
\end{enumerate}

\section{Monolithic versus hierarchical design}

monolithic VLA와 hierarchical VLA는 서로 배타적인 연구 방향이 아니다. task와 deployment condition에 따라 다른 점이 유리하다.

\begin{table}[h]
  \centering
  \caption{Monolithic VLA와 hierarchical VLA의 비교}
  \begin{tabularx}{\linewidth}{p{0.22\linewidth}X X}
    \toprule
    기준 & Monolithic VLA & Hierarchical VLA \\
    \midrule
    Short task & data가 충분하면 단순하고 빠름 & overhead가 생길 수 있음 \\
    Long horizon & compounding error와 memory 약점 & subgoal과 progress 관리에 유리 \\
    Safety & implicit behavior에 의존하기 쉬움 & external monitor와 constraint 삽입 가능 \\
    Debugging & failure 원인 분해가 어려움 & layer별 log와 failure taxonomy 가능 \\
    Latency & 한 model call로 끝날 수 있음 & reasoning call과 tool call 비용 존재 \\
    Adaptation & data scale이 크면 강력함 & robot-specific wrapper 교체가 쉬움 \\
    \bottomrule
  \end{tabularx}
  \label{tab:ch17-mono-hier}
\end{table}

공학적 결론은 단순하다. 짧고 반복적인 visuomotor task는 monolithic policy가 더 나을 수 있다. 그러나 multi-step manipulation, mobile manipulation, human-in-the-loop task, safety-critical task는 hierarchical structure가 더 많은 관측 가능성과 개입 지점을 제공한다.

\section{Safety gates}

hierarchical VLA에서 safety는 마지막에 붙이는 filter가 아니다. 각 layer에 gate가 있어야 한다.
\begin{equation}
  \act_t^{exec}
  =
  \mathcal{P}_{safe}
  (
  \act_t^{raw};
  \bel_t,
  c_t,
  \mathcal{S}_{robot},
  \mathcal{S}_{human}
  ).
  \label{eq:ch17-safety-projection}
\end{equation}
$\mathcal{P}_{safe}$는 raw action을 실행 가능한 safe action으로 projection하거나, projection이 불가능하면 halt 또는 ask로 전환한다.

\begin{table}[h]
  \centering
  \caption{계층별 safety gate}
  \begin{tabularx}{\linewidth}{p{0.2\linewidth}X X}
    \toprule
    Layer & Gate & 예시 \\
    \midrule
    Reasoning & instruction validity & 금지된 행동, ambiguous command, missing object \\
    Planning & feasibility & collision-free path, reachable grasp, precondition \\
    VLA policy & action sanity & saturation, invalid token, large pose jump \\
    Controller & physical limit & torque, velocity, joint limit, contact force \\
    Runtime monitor & emergency override & human proximity, unexpected contact, perception loss \\
    \bottomrule
  \end{tabularx}
  \label{tab:ch17-safety-gates}
\end{table}

이 구조는 model을 덜 신뢰하자는 말이 아니다. 어떤 model도 모든 physical constraint를 내부화했다는 evidence를 쉽게 제공할 수 없으므로, 검증 가능한 constraint는 system level에서 표현해야 한다는 뜻이다.

\section{Logging and evaluation}

hierarchical VLA를 평가하려면 final success만 기록해서는 안 된다. 각 layer의 decision이 남아야 한다.
\begin{equation}
  L_t
  =
  (
  \ell,
  z_t,
  p_t,
  \act_t^{raw},
  \act_t^{exec},
  s_t^{monitor},
  y_t
  ).
  \label{eq:ch17-log-record}
\end{equation}
이 log는 failure attribution에 필요하다. 실패가 reasoning error인지, object grounding error인지, VLA action error인지, controller tracking error인지, safety gate의 conservative rejection인지 구분해야 다음 iteration의 data collection이 가능하다.

평가 metric도 계층별로 나누어야 한다.
\begin{enumerate}
  \item reasoning accuracy: subgoal decomposition, object relation, constraint recognition
  \item progress accuracy: completion, failure, unsafe event detection
  \item action execution: tracking, smoothness, interruption latency
  \item recovery quality: replan success, ask/abort correctness, repeated failure reduction
  \item system cost: reasoning latency, tool call count, compute, human intervention
\end{enumerate}

\section{Design checklist}

hierarchical VLA를 설계할 때 다음 질문을 먼저 정해야 한다.
\begin{enumerate}
  \item 상위 reasoning layer는 매 step 호출되는가, event 기반으로 호출되는가?
  \item reasoning output은 자연어인가, typed command인가, hybrid인가?
  \item 하위 VLA는 action chunk, waypoint, low-level command 중 무엇을 출력하는가?
  \item progress detector는 visual model인가, state estimator인가, rule verifier인가?
  \item 실패 시 recovery option은 누가 선택하는가?
  \item safety gate는 action 전, controller 전, runtime 중 어디에 있는가?
  \item log는 layer별 decision과 rejected action을 보존하는가?
\end{enumerate}

\begin{casebox}{Case study: embodied reasoning layer를 과대해석하지 않기}
high-level reasoning layer는 task planning, spatial logic, success detection을 강화할 수 있지만, 그 자체가 low-level VLA policy나 verified controller는 아니다. Real VLA 관점에서는 reasoning output이 typed subgoal인지, action constraint인지, progress signal인지 분리하고, 하위 VLA와 controller가 무엇을 검증하는지 확인해야 한다.
\end{casebox}

\chapterwork
{SayCan/VoxPoser식 planner, OpenVLA/$\pi_0$식 action policy, embodied reasoning layer를 hierarchy 관점에서 비교하라 \cite{saycan2022,voxposer2023,openvla2024,pi02024}.}
{상위 reasoning layer가 실패했는지 하위 VLA action이 실패했는지 구분하려면 어떤 log가 필요한가?}
{long-horizon mobile manipulation task에 대해 reasoning layer, VLA policy, controller, monitor의 interface contract를 설계하라.}

\section{요약}

계층형 VLA는 고전 modular robotics로의 단순 회귀가 아니다. embodied reasoning model은 task semantics, scene relation, progress, failure, tool use를 다루고, 하위 VLA는 fluid한 visuomotor action generation을 담당하며, controller와 safety monitor는 physical execution을 검증한다. real deployment에서 중요한 것은 어느 layer가 더 지능적인가가 아니라, 어떤 정보가 어느 layer에서 결정되고 어떻게 검증되는가이다.

다음 장에서는 이 계층 구조가 실제 하드웨어와 시간 제약을 만날 때 어떤 문제가 생기는지 다룬다. VLA는 좋은 action을 예측하는 것만으로 충분하지 않다. 정해진 latency 안에, 제한된 compute와 sensor bandwidth 안에서, controller가 받아들일 수 있는 형태로 action을 내야 한다.
